\documentclass[11pt,a4paper]{ctexbook}
\usepackage[margin=2.2cm]{geometry}
\usepackage{hyperref}
\usepackage{xcolor}
\usepackage{longtable}
\usepackage{booktabs}
\usepackage{array}
\usepackage{enumitem}
\usepackage{fancyhdr}
\usepackage{titlesec}
\usepackage{tcolorbox}
\usepackage{listings}
\usepackage{fontspec}
\usepackage{graphicx}
\usepackage{tabularx}
\usepackage{makeidx}
\makeindex

\hypersetup{
  colorlinks=true,
  linkcolor=blue!60!black,
  urlcolor=blue!60!black,
  citecolor=blue!60!black,
  pdftitle={Claude Code Codex Strongest 完整指南},
  pdfauthor={Liu Jiarui},
  pdfsubject={Claude Code + Codex 一键安装、配置、使用与维护指南}
}

\definecolor{codebg}{RGB}{248,248,248}
\definecolor{codeframe}{RGB}{220,220,220}
\definecolor{warnbg}{RGB}{255,248,230}
\definecolor{infobg}{RGB}{235,247,255}
\definecolor{okbg}{RGB}{238,250,238}

\lstdefinestyle{cmd}{
  basicstyle=\ttfamily\small,
  backgroundcolor=\color{codebg},
  frame=single,
  rulecolor=\color{codeframe},
  breaklines=true,
  columns=fullflexible,
  keepspaces=true,
  showstringspaces=false
}
\lstset{style=cmd}

\newtcolorbox{infobox}{colback=infobg,colframe=blue!35!black,title=说明}
\newtcolorbox{warnbox}{colback=warnbg,colframe=orange!70!black,title=注意}
\newtcolorbox{successbox}{colback=okbg,colframe=green!45!black,title=验证通过}

\pagestyle{fancy}
\fancyhf{}
\lhead{Claude Code Codex Strongest}
\rhead{完整指南}
\cfoot{\thepage}

\title{\Huge Claude Code Codex Strongest\\[0.4em]\Large 完整安装、配置、使用、排错与维护指南}
\author{Liu Jiarui}
\date{2026年6月}

\begin{document}
\frontmatter
\maketitle
\tableofcontents

\chapter*{阅读说明}
\addcontentsline{toc}{chapter}{阅读说明}

本指南面向 \textbf{Claude Code Codex Strongest / claude-code-codex-strongest} 项目的使用者、维护者和二次开发者。它不仅说明如何双击一键安装，也解释安装器到底做了什么、每个配置目录放什么、Claude Code 与 Codex 的边界在哪里、如何验证安装是否成功，以及常见故障如何定位。

\begin{infobox}
项目定位：这是一个同时覆盖 \textbf{Claude Code} 与 \textbf{Codex CLI} 的开箱即用配置包。它把 Claude Code 的 skills、agents、commands、hooks、MCPs 和 Codex 的安全模板配置整合到一个跨平台一键安装器中。
\end{infobox}

\begin{warnbox}
安全边界：安装器不会写入 Codex token，不会复制 \texttt{auth.json}，不会把 Claude 登录态迁移给 Codex。Codex 登录必须由用户安装后手动执行 \texttt{codex login}。
\end{warnbox}

\mainmatter

\chapter{项目总览}

\section{项目名称与仓库}

项目展示名为 \textbf{Claude Code Codex Strongest}，仓库名为 \textbf{claude-code-codex-strongest}。GitHub 地址为：

\begin{lstlisting}
https://github.com/liujiarui0918/claude-code-codex-strongest
\end{lstlisting}

本项目原本是 Claude Code Strongest，后来扩展为同时安装和配置 Claude Code 与 Codex，因此命名为 Claude Code Codex Strongest。

\section{一句话解释}

这是一个面向 Windows 和 macOS 的双栈 AI 编程环境一键安装包：

\begin{itemize}[leftmargin=2em]
  \item 安装 \textbf{Claude Code CLI} 与 VS Code 插件 \texttt{anthropic.claude-code}；
  \item 安装 \textbf{Codex CLI} 与 VS Code 插件 \texttt{openai.chatgpt}；
  \item 部署 \texttt{\textasciitilde/.claude} 全套工作流配置；
  \item 部署 \texttt{\textasciitilde/.codex} 安全模板；
  \item 配置 8 个 MCP server；
  \item 安装 VS Code、Git、Node.js、uv、cc-switch 等基础工具。
\end{itemize}

\section{适合谁}

\begin{itemize}[leftmargin=2em]
  \item 从零开始使用 Claude Code / Codex 的新手；
  \item 想要完整工程实践配置而不是空白默认配置的中级用户；
  \item 需要统一 Windows / macOS 开发环境的团队；
  \item 想 fork 后定制自己 AI 编程工作流的高级用户。
\end{itemize}

\section{核心组成}

\begin{longtable}{p{0.22\textwidth}p{0.70\textwidth}}
\toprule
组成 & 作用 \\
\midrule
\texttt{install-windows.bat} & Windows 用户双击入口，下载 GitHub 最新 main 包并调用 PowerShell 安装器。 \\
\texttt{install-macos.command} & macOS 用户双击入口，下载 GitHub 最新 main 包并调用 Bash 安装器。 \\
\texttt{install/install-windows.ps1} & Windows 主安装器，负责安装依赖、部署 Claude/Codex 配置、验证安装。 \\
\texttt{install/install-macos.sh} & macOS 主安装器，负责 Homebrew、VS Code、CLI、扩展与配置部署。 \\
\texttt{\textasciitilde/.claude} & Claude Code 全局配置目录，包含 skills、agents、commands、hooks、docs 等。 \\
\texttt{\textasciitilde/.codex} & Codex 全局配置目录，包含 AGENTS.md、config.toml、.gitignore 等安全模板。 \\
\texttt{\textasciitilde/.claude.json} & Claude Code MCP server 配置文件。 \\
\texttt{.codex/} & 仓库内 Codex 安全模板，安装器会复制其中安全文件到用户 \texttt{\textasciitilde/.codex}。 \\
\bottomrule
\end{longtable}

\chapter{快速安装}

\section{Windows 一键安装}

推荐方式：从 Releases 页面下载 \texttt{install-windows.bat} 并双击运行。

\begin{lstlisting}[language={}]
https://github.com/liujiarui0918/claude-code-codex-strongest/releases/latest
\end{lstlisting}

安装器会自动完成：

\begin{enumerate}[leftmargin=2em]
  \item 安装或确认 VS Code；
  \item 创建/刷新 Windows 桌面 \texttt{Visual Studio Code.lnk}；
  \item 安装 Git、Node.js、uv；
  \item 安装 Claude Code CLI：\texttt{@anthropic-ai/claude-code}；
  \item 安装 Codex CLI：\texttt{@openai/codex}；
  \item 安装 VS Code 扩展：\texttt{anthropic.claude-code}、\texttt{openai.chatgpt}、中文语言包、Office Viewer；
  \item 安装 cc-switch；
  \item 部署 \texttt{\textasciitilde/.claude}；
  \item 部署 \texttt{\textasciitilde/.codex} 安全模板；
  \item 写入 \texttt{\textasciitilde/.claude.json} 的 8 个 MCP server；
  \item 运行安装验证。
\end{enumerate}

\subsection{Windows 一行命令}

\begin{lstlisting}[language={}]
irm https://raw.githubusercontent.com/liujiarui0918/claude-code-codex-strongest/main/install/bootstrap.ps1 | iex
\end{lstlisting}

\subsection{Windows 本地仓库运行}

\begin{lstlisting}[language={}]
git clone https://github.com/liujiarui0918/claude-code-codex-strongest.git
cd claude-code-codex-strongest
.\install\install-windows.ps1
\end{lstlisting}

\section{macOS 一键安装}

推荐方式：从 Releases 页面下载 \texttt{install-macos.command} 并双击运行。若 macOS 阻止运行，可右键选择“打开”。

\subsection{macOS 一行命令}

\begin{lstlisting}[language={}]
curl -fsSL https://raw.githubusercontent.com/liujiarui0918/claude-code-codex-strongest/main/install/bootstrap.sh | bash
\end{lstlisting}

\subsection{macOS 本地仓库运行}

\begin{lstlisting}[language={}]
git clone https://github.com/liujiarui0918/claude-code-codex-strongest.git
cd claude-code-codex-strongest
bash install/install-macos.sh
\end{lstlisting}

\chapter{安装器细节}

\section{Windows 安装器流程}

Windows 主安装器为 \texttt{install/install-windows.ps1}，要求 PowerShell 5.1 或更高版本。它的主要阶段如下：

\begin{longtable}{p{0.18\textwidth}p{0.74\textwidth}}
\toprule
阶段 & 说明 \\
\midrule
Phase 1 & 通过 winget / npm 安装 VS Code、Git、Node.js、uv、Claude Code CLI、Codex CLI、VS Code 扩展，并创建 VS Code 桌面快捷方式。 \\
Phase 2 & 部署仓库文件到 \texttt{\textasciitilde/.claude}。若目录被当前 Claude Code 会话占用，安装器会 fallback 到 robocopy 复制备份，而不是直接失败。 \\
Phase 3 & 从 \texttt{settings.template.json} 渲染 \texttt{settings.json}。默认不写 API key，除非用户显式传入参数。 \\
Phase 4 & 将 8 个 MCP server 写入 \texttt{\textasciitilde/.claude.json}。 \\
Codex & 将仓库 \texttt{.codex} 下的安全模板复制到 \texttt{\textasciitilde/.codex}，只复制 \texttt{AGENTS.md}、\texttt{config.toml}、\texttt{.gitignore}、\texttt{README.md}。 \\
Phase 5 & 验证 settings、docs、skills、agents、commands、hooks、Codex 文件、VS Code shortcut、VS Code 扩展、MCP server 数量。 \\
\bottomrule
\end{longtable}

\section{macOS 安装器流程}

macOS 主安装器为 \texttt{install/install-macos.sh}。它会：

\begin{itemize}[leftmargin=2em]
  \item 安装 Homebrew（若不存在）；
  \item 安装 VS Code、Git、Node.js、PowerShell 7；
  \item 安装 uv；
  \item 安装 Claude Code CLI 与 Codex CLI；
  \item 确保 \texttt{code} 命令可用；
  \item 安装四个 VS Code 扩展；
  \item 部署 \texttt{\textasciitilde/.claude}；
  \item 渲染 Claude Code settings；
  \item 写入 MCP server；
  \item 调用 \texttt{deploy\_codex\_config} 部署 Codex 模板；
  \item 验证安装。
\end{itemize}

\section{Bootstrap 脚本}

\texttt{install/bootstrap.ps1} 与 \texttt{install/bootstrap.sh} 是“无需 clone”的远程入口。它们会下载 GitHub main 分支压缩包，解压 \texttt{claude-code-codex-strongest-*} 目录，再调用平台主安装器。

\subsection{固定版本或分支}

两个 bootstrap 都支持通过 \texttt{CCS\_REF} 固定版本、分支或 tag。若值形如版本号或 tag，脚本会走 GitHub tags；否则走 heads。

Windows PowerShell：

\begin{lstlisting}[language={}]
$env:CCS_REF = 'v1.0.0'
irm https://raw.githubusercontent.com/liujiarui0918/claude-code-codex-strongest/main/install/bootstrap.ps1 | iex
\end{lstlisting}

macOS / Linux：

\begin{lstlisting}[language={}]
CCS_REF=v1.0.0 curl -fsSL https://raw.githubusercontent.com/liujiarui0918/claude-code-codex-strongest/main/install/bootstrap.sh | bash
\end{lstlisting}

\chapter{安装内容清单}

\section{基础工具}

\begin{longtable}{p{0.25\textwidth}p{0.32\textwidth}p{0.32\textwidth}}
\toprule
工具 & Windows 渠道 & macOS 渠道 \\
\midrule
VS Code & winget \texttt{Microsoft.VisualStudioCode} & brew cask \texttt{visual-studio-code} \\
Git & winget \texttt{Git.Git} & brew \texttt{git} \\
Node.js LTS & winget \texttt{OpenJS.NodeJS} & brew \texttt{node} \\
uv & winget \texttt{astral-sh.uv} & 官方安装脚本 \\
PowerShell & Windows 自带 5.1 可用 & brew cask \texttt{powershell} \\
cc-switch & winget \texttt{farion1231.CC-Switch} & brew cask \texttt{cc-switch} \\
\bottomrule
\end{longtable}

\section{CLI 与 VS Code 扩展}

\begin{longtable}{p{0.26\textwidth}p{0.62\textwidth}}
\toprule
项目 & 安装内容 \\
\midrule
Claude Code CLI & npm 包 \texttt{@anthropic-ai/claude-code}，命令 \texttt{claude}。 \\
Claude Code VS Code 扩展 & 扩展 ID \texttt{anthropic.claude-code}。 \\
Codex CLI & npm 包 \texttt{@openai/codex}，命令 \texttt{codex}。 \\
Codex / ChatGPT VS Code 扩展 & 扩展 ID \texttt{openai.chatgpt}。 \\
中文语言包 & 扩展 ID \texttt{MS-CEINTL.vscode-language-pack-zh-hans}。 \\
Office Viewer & 扩展 ID \texttt{cweijan.vscode-office}。 \\
\bottomrule
\end{longtable}



\chapter{为什么是 Strongest：完整配置资产逐项拆解}

\begin{infobox}
本章是全书亮点。Claude Code Codex Strongest 的强，不是因为名字里有 Strongest，而是因为它把 Claude Code 与 Codex 两套 CLI、VS Code 插件、cc-switch 模型路由、CLAUDE.md 全局规则、docs 方法文档、33 个 skills、22 个 agents、25 个 commands、16 个 hook 脚本、settings.json 生命周期配置、MCP servers、安全边界和真实验证流程整合为一套可安装、可复用、可审计的工程系统。
\end{infobox}

\section{Strongest 的判定标准}

一个真正强的 AI 编程配置，不应该只回答“能不能聊天”，而要回答五个工程问题：第一，普通用户能否一键装好；第二，复杂任务能否被拆成计划、实现、审查、验证；第三，危险操作和密钥泄漏能否提前拦住；第四，模型、工具、文档、浏览器、Git 能否接入同一条工作流；第五，产物能否被真实命令证明可用。

本项目的答案是：用一键安装器负责环境，用 \texttt{CLAUDE.md} 和 docs 负责规则，用 skills 负责方法，用 agents 负责角色分工，用 commands 负责入口，用 hooks 负责安全和自动化，用 settings.json 负责生命周期注册，用 MCP 负责外部能力，用 Codex 的 \texttt{AGENTS.md} 和 \texttt{config.toml} 负责第二套 CLI 的隔离配置。

\section{CLAUDE.md：全局入口，不是普通 README}

\texttt{CLAUDE.md} 是 Claude Code 每次会话自动加载的高优先级指令入口。它本身很短，但通过模块引用把真正的规则拆到四个文档：

\begin{longtable}{p{0.22\textwidth}p{0.66\textwidth}}
\toprule
文件 & 作用 \\
\midrule
CLAUDE.md & 会话总入口。声明全局规则、模块引用、已安装 skills/agents/commands/MCPs/hooks、备份和 doctor 命令。 \\
docs/environment.md & 定义 OS、Shell、PowerShell 兼容规则、路径约定、不要用 Unix-only 命令等环境规则。 \\
docs/workflow.md & 定义沟通风格、非平凡任务先 plan、Edit 前必须 Read、验证要求、研究搜索策略。 \\
docs/tools.md & 定义 skills、subagents、MCP、cron 的使用边界，强调优先用工具而不是重头实现。 \\
docs/safety.md & 定义 Git、安全、destructive ops、反模式，明确 reset/force push/删除目录/改系统配置前必须确认。 \\
\bottomrule
\end{longtable}

这种拆分的强点在于：\texttt{CLAUDE.md} 不臃肿，规则又足够完整；每个模块可以单独维护，安装器只要部署整套目录，就能让新会话自动继承工作方法。


\section{小白图文流程：从中转站到 cc-switch}

下面九张图来自 \texttt{docs/images/}，对应 \texttt{guide-for-beginners.md} 的完整上手流程。它们被合并进 PDF，是为了让没有任何基础的用户也能按图操作。

\subsection{中转站控制台：URL、Key、模型名}

图 \ref{fig:relay-api-info} 展示中转站控制台右侧 API 信息区域，用户需要复制一个接口 URL。注意 URL 末尾不要加 \texttt{/}。

\begin{figure}[htbp]
\centering
\includegraphics[width=0.82\textwidth]{images/1.png}
\caption{中转站控制台：右侧 API 信息区域显示多个可用域名。}
\label{fig:relay-api-info}
\end{figure}

图 \ref{fig:relay-token} 展示令牌管理页面，用户从这里复制 API Key。Key 只能放在 cc-switch 或脚本参数里，不能提交到 Git，也不能写进文档。

\begin{figure}[htbp]
\centering
\includegraphics[width=0.82\textwidth]{images/2.png}
\caption{令牌管理页：复制 API Key。}
\label{fig:relay-token}
\end{figure}

图 \ref{fig:model-id} 展示模型广场里模型详情页右上角的模型 ID。不同中转站模型名可能不同，必须以控制台显示为准。

\begin{figure}[htbp]
\centering
\includegraphics[width=0.82\textwidth]{images/3.png}
\caption{模型详情：复制模型 ID。}
\label{fig:model-id}
\end{figure}

图 \ref{fig:model-price} 展示分组价格。相同模型在不同分组可能价格差很多，选择令牌前要确认分组和倍率。

\begin{figure}[htbp]
\centering
\includegraphics[width=0.82\textwidth]{images/4.png}
\caption{模型详情：查看分组价格，选择成本合适的分组。}
\label{fig:model-price}
\end{figure}

\subsection{cc-switch：添加 provider、模型映射、路由代理}

图 \ref{fig:ccswitch-provider} 展示 cc-switch 添加 provider 的界面。这里填写供应商名称、API Key 和请求地址。

\begin{figure}[htbp]
\centering
\includegraphics[width=0.82\textwidth]{images/5.png}
\caption{cc-switch provider 编辑页：填写 API Key 和请求地址。}
\label{fig:ccswitch-provider}
\end{figure}

图 \ref{fig:ccswitch-models} 展示获取模型列表和模型映射。推荐把 Opus 映射到最强模型，Sonnet 映射到日常模型，Haiku 映射到低成本执行模型。

\begin{figure}[htbp]
\centering
\includegraphics[width=0.82\textwidth]{images/6.png}
\caption{cc-switch 获取模型列表：为 Opus / Sonnet / Haiku 配置映射。}
\label{fig:ccswitch-models}
\end{figure}

图 \ref{fig:openai-endpoint} 展示模型端点类型。如果显示 \texttt{openai: /v1/chat/completions}，Claude Code 侧需要开启 cc-switch 的 Claude 路由转换。

\begin{figure}[htbp]
\centering
\includegraphics[width=0.82\textwidth]{images/7.png}
\caption{模型广场：端点显示 openai，说明需要路由转换。}
\label{fig:openai-endpoint}
\end{figure}

图 \ref{fig:ccswitch-route} 展示路由设置页。需要打开路由总开关，并启用 Claude 路由。

\begin{figure}[htbp]
\centering
\includegraphics[width=0.82\textwidth]{images/8.png}
\caption{cc-switch 路由设置：打开总开关和 Claude 路由。}
\label{fig:ccswitch-route}
\end{figure}

图 \ref{fig:ccswitch-proxy-green} 展示主界面代理图标变绿，表示路由代理已启用。

\begin{figure}[htbp]
\centering
\includegraphics[width=0.82\textwidth]{images/9.png}
\caption{cc-switch 主界面：代理图标变绿表示路由已启用。}
\label{fig:ccswitch-proxy-green}
\end{figure}


\section{33 个 Skills：把专家方法固化成可触发流程}

Skills 是本项目最核心的“方法论层”。它们不是简单提示词，而是给不同任务类型准备的操作规程。下面逐个列出：

\begin{longtable}{p{0.24\textwidth}p{0.38\textwidth}p{0.28\textwidth}}
\toprule
Skill & 作用 & 典型场景 \\
\midrule
mcp-builder & 构建新的 MCP server，覆盖 stdio/HTTP、tool schema、错误处理。 & 当要给 Claude/Codex 增加外部工具能力时使用。 \\
document-handling & 处理 docx/xlsx/pptx/pdf 等二进制文档。 & 生成或读取 Office/PDF 文件，避免用 cat/grep 读二进制。 \\
lean-context & 控制上下文体积，只读取必要内容。 & 长任务、长文件、多轮会话中防止上下文污染。 \\
using-superpowers & 会话开始时选择正确 skill。 & 进入新任务时先判断该用哪个工作方法。 \\
brainstorming & 设计前澄清需求、比较 2--3 个方案。 & 需求不清、架构选择、产品方案讨论。 \\
writing-plans & 把方案拆成可验证的执行任务。 & 准备落地复杂功能或修复前写 PLAN。 \\
executing-plans & 按 1--3 个任务一批执行、验证、汇报。 & 已经有计划文档，需要稳步落地。 \\
subagent-driven-development & 复杂任务每个任务派 fresh subagent，并做 spec review。 & 计划很长、文件很多、主上下文会被污染。 \\
dispatching-parallel-agents & 多个独立任务并行分派。 & 搜索、审查、验证等可以并行完成的工作。 \\
test-driven-development & 严格 RED--GREEN--REFACTOR。 & 明确 bug 复现或明确 API/行为规格。 \\
systematic-debugging & 复现、隔离、修复、验证四阶段调试。 & 任何真实 bug、安装失败、CLI 异常。 \\
verification-before-completion & 完成前必须跑真实验证并读输出。 & 准备声明 done、提交、发布前。 \\
requesting-code-review & 交付前自检和请求代码审查。 & 完成改动但还没交付/commit 前。 \\
receiving-code-review & 处理审查反馈，不盲从也不防御。 & 收到 reviewer 或用户指出问题后。 \\
using-git-worktrees & 用 git worktree 隔离 feature。 & 需要在不污染 main 的情况下做大改动。 \\
finishing-a-development-branch & 收尾分支并给出 merge/PR/keep/discard 选项。 & 功能完成、测试通过后准备收口。 \\
writing-skills & 编写新 skill 的模板和反模式。 & 要把某套工作方法沉淀成 skill。 \\
condition-based-waiting & 异步等待必须等条件，不盲目 sleep。 & 服务启动、CI、后台任务、文件生成。 \\
root-cause-tracing & 沿调用链追到根因，不停在症状。 & 调试和修复阶段防止补丁式修复。 \\
defense-in-depth & 关键不变量需要多层防护。 & 安全、数据删除、权限、凭证保护。 \\
using-context7 & 查实时库/框架/SDK 文档。 & 涉及库 API、配置、版本迁移时。 \\
using-deepwiki & 理解公开 GitHub repo 的结构和设计。 & 研究外部仓库、框架、工具实现。 \\
using-sequential-thinking & 复杂多步推理和架构选择。 & 需要分支、修正、逐步推理的问题。 \\
using-playwright & 浏览器自动化、截图和 UI 验证。 & 前端/UI/网页抓取/真实交互验证。 \\
using-mcp-memory & 使用知识图谱 memory MCP。 & 需要实体、关系、观察记录的长期知识。 \\
using-cron & 提醒、定时、循环任务。 & 每 N 分钟检查、每日运行、定时提醒。 \\
using-task-list & 判断是否需要任务列表并维护任务。 & 3+ 步可追踪工作或多人/多 agent 协作。 \\
using-claude-api & Claude API / Anthropic SDK 参考。 & 写 Anthropic SDK、模型、tool use、缓存、流式。 \\
extracting-from-images & 从截图/图片提取文字、布局、意图。 & 用户贴图、UI 截图、错误截图。 \\
prompt-engineering & 编写和改进 LLM prompt。 & 需要结构化输出、few-shot、评测 prompt。 \\
incremental-context-building & 进入陌生大代码库时分层建立理解。 & 新项目 onboarding、大型代码库探索。 \\
skill-creator & 创建新 skill 的标准化流程。 & 把成熟流程产品化成可复用 skill。 \\
using-memory & 使用 auto-memory 系统，判断什么该记。 & 用户要求记住偏好或项目事实时。 \\
\bottomrule
\end{longtable}

\section{22 个 Agents：把复杂任务交给专门角色}

Agents 是“角色层”。当任务需要大量上下文、专项知识或并行处理时，主线程不应该自己吞下所有文件，而应把任务分给专门 agent。

\begin{longtable}{p{0.26\textwidth}p{0.62\textwidth}}
\toprule
Agent & 职责 \\
\midrule
planner & 规划非平凡实现，输出步骤、关键文件、风险。 \\
code-reviewer & 代码质量、可维护性、正确性审查，只读。 \\
security-reviewer & OWASP、secrets、注入、鉴权鉴权审查。 \\
tdd-guide & 用 TDD 驱动新功能或 bug 修复。 \\
debugger & 系统性复现、隔离、修复、验证 bug。 \\
test-runner & 运行测试并解释失败。 \\
refactor-assistant & 安全重构，保持行为不变。 \\
doc-writer & README、API 文档、变更日志写作。 \\
api-designer & REST/GraphQL/gRPC API schema 设计。 \\
performance-analyst & 先 profiling 后优化，定位性能热点。 \\
pr-reviewer & PR 描述与 diff review。 \\
research-agent & 多步研究、资料综合。 \\
explore & 快速只读代码库探索、定位符号和引用。 \\
build-fixer & 修复构建、类型、lint、工具链错误。 \\
typescript-reviewer & TypeScript 专项审查，关注 any、断言、泛型、async。 \\
python-reviewer & Python 专项审查，关注可变默认值、async、类型。 \\
git-operator & 高级 Git 操作，如 rebase、cherry-pick、reflog 恢复。 \\
regex-expert & 编写、解释、调试正则。 \\
sql-expert & SQL 查询、schema、索引、迁移和 explain。 \\
prompt-engineer & Prompt 设计、结构化输出、反模式。 \\
migration-assistant & 语言/框架迁移，如 Python、React、Django。 \\
accessibility-auditor & WCAG、ARIA、键盘、屏幕阅读器、对比度审查。 \\
\bottomrule
\end{longtable}

\section{25 个 Commands：把工作流变成一条入口}

Commands 是“入口层”。新手不用背完整 prompt，只要知道自己要规划、调试、审查、验证，就能通过 slash command 进入正确流程。

\begin{longtable}{p{0.24\textwidth}p{0.64\textwidth}}
\toprule
Command & 用途 \\
\midrule
/plan & 进入规划模式，先 brainstorm 再写可执行计划。 \\
/tdd & 用严格 TDD 实现功能或修 bug。 \\
/verify & 运行真实命令验证 claim。 \\
/review & 审查当前 diff 或指定路径。 \\
/security-scan & 安全扫描：OWASP、secrets、注入、鉴权。 \\
/debug & 系统性调试：复现、隔离、修复、验证。 \\
/refactor & 安全重构，先看测试覆盖。 \\
/worktree & 创建隔离 git worktree。 \\
/deep-research & 多源深度研究并综合报告。 \\
/ship & 分支收尾：验证后选择 merge/PR/keep/discard。 \\
/onboard & 新代码库 onboarding，识别结构、栈和关键文件。 \\
/think & 用 sequential-thinking 做复杂推理。 \\
/explain & 解释代码、错误或概念。 \\
/gen-tests & 为文件或函数生成测试。 \\
/changelog & 根据 git 历史生成 Keep a Changelog 条目。 \\
/eli5 & 用类比给小白解释技术概念。 \\
/regex & 按需求编写正则和测试用例。 \\
/cleanup & 清理当前 diff 中死代码、未用 import、TODO、console.log。 \\
/deps-audit & 审计依赖过期、漏洞、未使用。 \\
/skills & 列出可用 skills。 \\
/memory-search & 搜索 auto-memory。 \\
/mcp-status & 检查 MCP 配置状态。 \\
/backup & 备份 Claude Code 配置为 zip。 \\
/agents & 列出可用 agents。 \\
/doctor & 诊断 Claude Code 配置健康。 \\
\bottomrule
\end{longtable}

\section{Hooks：把安全、质量和上下文自动化}

Hooks 是“安全网”。模型可能忘记规则，但 hook 在工具调用前后执行，能把关键边界从“靠自觉”变成“靠机制”。仓库中 hook 脚本如下：

\begin{longtable}{p{0.26\textwidth}p{0.24\textwidth}p{0.40\textwidth}}
\toprule
Hook & 注册/用途位置 & 作用 \\
\midrule
block-dangerous.ps1 & PreToolUse Bash/PowerShell & 拦截 git reset --hard、force push、--no-verify、系统关机/格式化等危险命令。 \\
protect-secrets.ps1 & PreToolUse Read/Edit/Write/Bash/PowerShell & 保护 .env、SSH private key、AWS credentials、auth.json、token 等敏感文件。 \\
protect-linter-configs.ps1 & PreToolUse Edit/Write/MultiEdit & 防止为了通过检查而修改 ESLint、Prettier、tsconfig 等质量门配置。 \\
auto-format.ps1 & PostToolUse Edit/Write & 编辑后自动格式化，减少风格噪音。 \\
typecheck.ps1 & PostToolUse Edit/Write async & 编辑后异步类型检查，尽早发现破坏。 \\
command-log.ps1 & PostToolUse Bash/PowerShell & 记录命令审计日志，方便追溯。 \\
session-start-context.ps1 & SessionStart & 注入 git 状态、最近提交、可用 skills/agents/MCPs。 \\
smart-context.ps1 & UserPromptSubmit & 根据用户 prompt 提示合适 skill 和注意事项。 \\
notify.ps1 & Notification & 桌面通知。 \\
stop-sound.ps1 & Stop & 任务完成提示音。 \\
precompact-backup.ps1 & PreCompact & 压缩上下文前备份 transcript。 \\
subagent-stop.ps1 & SubagentStop & 记录 subagent 结束事件。 \\
statusline.ps1 & statusLine & 每 10 秒刷新状态行，显示项目、分支、模型等上下文。 \\
backup-config.ps1 & 命令辅助 & 供 /backup 调用，备份配置。 \\
mcp-doctor.ps1 & 诊断辅助 & 检查 MCP 配置和依赖。 \\
post-compact-restore.ps1 & 压缩辅助 & 压缩后恢复或提示上下文。 \\
\bottomrule
\end{longtable}

\section{settings.json：生命周期编排核心}

\texttt{settings.template.json} 是 Claude Code 侧的生命周期编排模板。安装器不会简单复制固定 \texttt{settings.json}，而是渲染占位符：\texttt{\{\{CLAUDE\_HOME\}\}}、\texttt{\{\{ANTHROPIC\_AUTH\_TOKEN\}\}}、\texttt{\{\{ANTHROPIC\_BASE\_URL\}\}}、\texttt{\{\{ANTHROPIC\_MODEL\}\}}、\texttt{\{\{TIMEZONE\}\}} 等。

\subsection{环境变量设计}

\begin{longtable}{p{0.34\textwidth}p{0.54\textwidth}}
\toprule
变量 & 说明 \\
\midrule
ANTHROPIC\_AUTH\_TOKEN & 只有用户显式传入 token 时才写入；默认删除，避免安装器保存凭证。 \\
ANTHROPIC\_BASE\_URL & 只有用户显式指定 Base URL 时写入；通常推荐由 cc-switch 管理。 \\
ANTHROPIC\_MODEL & 可选模型钉住，用于中转站或脚本化安装。 \\
ANTHROPIC\_DEFAULT\_HAIKU\_MODEL & 与 ANTHROPIC\_MODEL 同步，兼容后台轻量模型调用。 \\
EDITOR & 默认 code --wait，让编辑操作进入 VS Code。 \\
CLAUDE\_HOOKS\_LOG\_DIR & 指向 \textasciitilde{}/.claude/logs，供 hooks 写审计和诊断日志。 \\
CLAUDE\_DISABLED\_HOOKS & 保留禁用 hook 的开关，便于临时排障。 \\
\bottomrule
\end{longtable}

\subsection{statusLine 与 permissions}

\texttt{statusLine} 每 10 秒执行 \texttt{hooks/statusline.ps1}。它让会话底部能显示项目、分支、模型或其他上下文。\texttt{permissions.deny} 默认拒绝高危 PowerShell 命令：\texttt{Format-Volume}、\texttt{Stop-Computer}、\texttt{Restart-Computer}、\texttt{Clear-RecycleBin}、\texttt{Remove-PSDrive}。这属于防御纵深：即使模型误判，权限层和 hook 层也会先拦一次。

\subsection{生命周期注册矩阵}

\begin{longtable}{p{0.18\textwidth}p{0.22\textwidth}p{0.24\textwidth}p{0.26\textwidth}}
\toprule
事件 & Matcher & 脚本 & 效果 \\
\midrule
PreToolUse & Bash|PowerShell & block-dangerous.ps1 & 命令执行前拦截危险 shell/Git/系统操作。 \\
PreToolUse & Read|Edit|Write|Bash|PowerShell & protect-secrets.ps1 & 读取或修改前保护密钥类文件。 \\
PreToolUse & Edit|Write|MultiEdit & protect-linter-configs.ps1 & 防止通过改 linter 配置绕过质量检查。 \\
PostToolUse & Edit|Write & auto-format.ps1 & 编辑后自动格式化。 \\
PostToolUse & Edit|Write & typecheck.ps1 & 编辑后异步类型检查。 \\
PostToolUse & Bash|PowerShell & command-log.ps1 & 记录命令审计。 \\
SessionStart & - & session-start-context.ps1 & 会话开始注入 git 状态和项目上下文。 \\
UserPromptSubmit & - & smart-context.ps1 & 根据用户 prompt 提醒合适 skill。 \\
Notification & - & notify.ps1 & 桌面通知。 \\
Stop & - & stop-sound.ps1 & 结束提示音。 \\
PreCompact & - & precompact-backup.ps1 & 上下文压缩前备份 transcript。 \\
SubagentStop & - & subagent-stop.ps1 & 记录 subagent 停止。 \\
\bottomrule
\end{longtable}

\section{为什么这些组合起来才是 Strongest}

单独的 \texttt{CLAUDE.md} 只是规则；单独的 skill 只是方法；单独的 agent 只是角色；单独的 hook 只是拦截器；单独的 MCP 只是工具。Claude Code Codex Strongest 的强点，是把它们串成闭环：用户提出需求后，smart-context 提醒方法，skills 定义流程，commands 提供入口，agents 分担复杂任务，hooks 防止越界，MCP 补充外部能力，settings.json 负责注册，doctor/verify 给出证据，Codex 作为独立第二栈补充执行与对照能力。

这就是第五章要论证的重点：Strongest 不是“大而全”的堆料，而是“每一层都有职责，每一层都能被安装器部署，每一层都能被验证”的工程化系统。

\chapter{Codex 配置}

\section{目录与模板}

Codex 的运行目录是 \texttt{\textasciitilde/.codex}。仓库内的 \texttt{.codex/} 是安全模板目录，安装器只复制这些文件：

\begin{itemize}[leftmargin=2em]
  \item \texttt{AGENTS.md}：Codex 可读的工作规则入口；
  \item \texttt{config.toml}：Codex 基础配置；
  \item \texttt{.gitignore}：防止 auth、token、runtime 状态被提交；
  \item \texttt{README.md}：模板说明。
\end{itemize}

\section{不会复制什么}

\begin{warnbox}
安装器不会复制或生成 \texttt{auth.json}、\texttt{.credentials.json}、token、key、sessions、logs、cache、runtime 等状态。用户必须通过 \texttt{codex login} 自己建立 Codex 登录态。
\end{warnbox}

\section{Codex 登录}

安装结束后，运行：

\begin{lstlisting}[language={}]
codex login
\end{lstlisting}

登录后可以验证：

\begin{lstlisting}[language={}]
codex --version
codex doctor --summary --ascii --no-color
\end{lstlisting}

在一次真实验证中，本机输出了：

\begin{lstlisting}[language={}]
codex-cli 0.139.0
Codex Doctor v0.139.0 ...
15 ok | 1 idle | 2 notes | 2 warn | 0 fail degraded
\end{lstlisting}

其中 WebSocket timeout warning 不等于配置失败；doctor 明确显示 HTTPS fallback 可能仍可用，且 exit code 为 0。

\chapter{MCP Servers}

安装器会把 8 个 MCP server 写入 \texttt{\textasciitilde/.claude.json}：

\begin{longtable}{p{0.28\textwidth}p{0.58\textwidth}}
\toprule
MCP & 用途 \\
\midrule
\texttt{sequential-thinking} & 多步推理与复杂决策。 \\
\texttt{context7} & 实时库文档查询。 \\
\texttt{deepwiki} & 公开 GitHub repo 文档问答。 \\
\texttt{playwright} & 浏览器自动化、截图、UI 验证。 \\
\texttt{memory} & 知识图谱式记忆。 \\
\texttt{fetch} & HTTP 抓取。 \\
\texttt{time} & 时区转换，默认 Asia/Shanghai。 \\
\texttt{git-mcp} & Git 操作。 \\
\bottomrule
\end{longtable}

\section{Windows 与 macOS 的命令差异}

Windows 模板将 stdio MCP 命令包在 \texttt{cmd /c} 后面，例如 \texttt{cmd /c npx -y ...} 或 \texttt{cmd /c uvx ...}。这样可以提升 Windows 下 npm、uvx、PATH 解析的稳定性。macOS 模板则直接调用 \texttt{npx} 或 \texttt{uvx}。

\begin{longtable}{p{0.22\textwidth}p{0.34\textwidth}p{0.34\textwidth}}
\toprule
MCP & Windows & macOS \\
\midrule
\texttt{sequential-thinking} & \texttt{cmd /c npx -y @modelcontextprotocol/server-sequential-thinking} & \texttt{npx -y @modelcontextprotocol/server-sequential-thinking} \\
\texttt{playwright} & \texttt{cmd /c npx -y @playwright/mcp@latest --headless} & \texttt{npx -y @playwright/mcp@latest --headless} \\
\texttt{memory} & \texttt{cmd /c npx -y @modelcontextprotocol/server-memory} & \texttt{npx -y @modelcontextprotocol/server-memory} \\
\texttt{fetch} & \texttt{cmd /c uvx mcp-server-fetch} & \texttt{uvx mcp-server-fetch} \\
\texttt{time} & \texttt{cmd /c uvx mcp-server-time --local-timezone <TZ>} & \texttt{uvx mcp-server-time --local-timezone <TZ>} \\
\texttt{git-mcp} & \texttt{cmd /c uvx mcp-server-git} & \texttt{uvx mcp-server-git} \\
\bottomrule
\end{longtable}

\section{写入策略}

安装器会把模板里的 \texttt{\{\{TIMEZONE\}\}} 替换为用户传入的 IANA 时区，默认 \texttt{Asia/Shanghai}。随后覆盖写入 \texttt{\textasciitilde/.claude.json} 的 \texttt{mcpServers} 属性；如果 \texttt{hasCompletedOnboarding} 缺失，则设为 \texttt{true}，避免 fresh/reset 用户重复触发 onboarding。

\chapter{API Key 与模型配置}

\section{Claude Code 侧：cc-switch}

本项目取消了安装器中的 API key 弹窗。安装完成后，cc-switch 会打开，用户在其中添加 provider：

\begin{enumerate}[leftmargin=2em]
  \item 点击 Add Provider；
  \item 输入 API key 与 Base URL；
  \item 选择 Claude/Anthropic 预设，或使用中转站配置；
  \item 点击 Enable；
  \item cc-switch 写入 Claude Code 所需配置。
\end{enumerate}

\section{API 资源链接}

配置 provider 时可以参考以下入口：

\begin{itemize}[leftmargin=2em]
  \item 全网最低价 API 价格表：\url{https://docs.qq.com/sheet/DWFhhV3pMVnJGZXpO?tab=000001}
  \item Claude 公益站（每天 25 美元额度）：\url{https://anyrouter.top/register?aff=rQEI}
\end{itemize}

\section{Codex 侧：codex login}

Codex 登录不通过安装器写入，必须手动执行：

\begin{lstlisting}[language={}]
codex login
\end{lstlisting}

\section{模型名与中转站}

对于 Claude Code，安装器支持：

\begin{itemize}[leftmargin=2em]
  \item \texttt{-ApiToken} / \texttt{--token}：脚本化写入 token；
  \item \texttt{-BaseUrl} / \texttt{--url}：指定 API base URL；
  \item \texttt{-Model} / \texttt{--model}：同时设置 \texttt{ANTHROPIC\_MODEL} 与 \texttt{ANTHROPIC\_DEFAULT\_HAIKU\_MODEL}；
  \item \texttt{-NoCcSwitch} / \texttt{--no-cc-switch}：跳过 cc-switch。
\end{itemize}

\chapter{常用操作}

\section{验证安装}

Windows 下可运行：

\begin{lstlisting}[language={}]
claude --version
codex --version
codex doctor --summary --ascii --no-color
code --list-extensions | findstr /i "anthropic.claude-code openai.chatgpt"
\end{lstlisting}

Git Bash / macOS 下可运行：

\begin{lstlisting}[language={}]
claude --version
codex --version
codex doctor --summary --ascii --no-color
code --list-extensions | grep -Ei '^(anthropic\.claude-code|openai\.chatgpt)$'
\end{lstlisting}

\section{Reset 干净重装}

Windows：

\begin{lstlisting}[language={}]
.\install\install-windows.ps1 -Reset
\end{lstlisting}

macOS：

\begin{lstlisting}[language={}]
bash install/install-macos.sh --reset
\end{lstlisting}

\section{非交互安装}

Windows：

\begin{lstlisting}[language={}]
.\install\install-windows.ps1 -NonInteractive -Force
\end{lstlisting}

macOS：

\begin{lstlisting}[language={}]
bash install/install-macos.sh --non-interactive --force
\end{lstlisting}

\chapter{排错指南}

\section{GitHub 下载失败}

现象：bootstrap 下载 zip/tar.gz 失败。原因通常是网络、GitHub 访问、代理或证书问题。

处理：

\begin{itemize}[leftmargin=2em]
  \item 确认 VPN / 代理；
  \item 直接访问 Releases 页面下载入口脚本；
  \item clone 仓库后本地运行安装器。
\end{itemize}

\section{winget msstore 证书错误}

曾观察到：

\begin{lstlisting}[language={}]
0x8a15005e : 服务器证书与任何预期值都不匹配
\end{lstlisting}

若包在 winget 源存在，可显式指定：

\begin{lstlisting}[language={}]
winget install -e --id <Package.Id> --source winget
\end{lstlisting}

\section{VS Code 扩展安装出现 Node DeprecationWarning}

VS Code CLI 可能输出 \texttt{DEP0169} warning。Windows 安装器已修复：不再把 stderr warning 当成安装失败，而是检查 \texttt{\$LASTEXITCODE}。

\section{正在 Claude Code 内运行安装器导致 .claude 被占用}

现象：备份 \texttt{\textasciitilde/.claude} 时提示“另一个程序正在使用此文件”。

当前 Windows 安装器已修复：

\begin{itemize}[leftmargin=2em]
  \item 先尝试 Move-Item 整目录备份；
  \item 若失败，使用 robocopy 复制备份；
  \item 不删除正在使用的目录；
  \item 继续 overlay 部署。
\end{itemize}

\section{Codex doctor WebSocket warning}

若看到：

\begin{lstlisting}[language={}]
Responses WebSocket timed out; HTTPS fallback may still work
\end{lstlisting}

这通常不是安装失败。只要 doctor exit code 为 0，且没有 fail 项，就可以继续使用；若实际请求失败，再检查代理、VPN、防火墙、DNS、证书和 WebSocket 策略。

\section{codex 命令找不到}

处理顺序：

\begin{enumerate}[leftmargin=2em]
  \item 重新打开终端，让 PATH 刷新；
  \item 检查 npm 全局 bin 是否在 PATH；
  \item 运行 \texttt{npm install -g @openai/codex}；
  \item 再运行 \texttt{codex --version}。
\end{enumerate}

\chapter{安全模型}

\section{不迁移凭证}

项目的核心安全原则：安装器可以部署配置，但不能替用户迁移或生成凭证。

禁止复制或提交：

\begin{itemize}[leftmargin=2em]
  \item \texttt{auth.json}
  \item \texttt{.credentials.json}
  \item \texttt{*.token}
  \item \texttt{*.key}
  \item \texttt{.env} 与 \texttt{.env.*}
  \item sessions、history、logs、tmp、cache、runtime、backups
\end{itemize}

\section{Git 安全}

默认规则：

\begin{itemize}[leftmargin=2em]
  \item 不使用 \texttt{--no-verify}；
  \item 不 force push，除非用户明确要求；
  \item 不提交 secrets；
  \item destructive operations 前必须确认。
\end{itemize}

\chapter{维护与发布}

\section{仓库重命名}

项目仓库已重命名为：

\begin{lstlisting}[language={}]
liujiarui0918/claude-code-codex-strongest
\end{lstlisting}

本地 remote 应为：

\begin{lstlisting}[language={}]
origin https://github.com/liujiarui0918/claude-code-codex-strongest.git
\end{lstlisting}

\section{提交前检查}

建议每次发布前运行：

\begin{lstlisting}[language={}]
git status --short
powershell -NoProfile -ExecutionPolicy Bypass -Command "$p='install/install-windows.ps1'; $tokens=$null; $errors=$null; [System.Management.Automation.Language.Parser]::ParseFile($p,[ref]$tokens,[ref]$errors) | Out-Null; if($errors.Count){ exit 1 }"
bash -n install/install-macos.sh
bash -n install-macos.command
\end{lstlisting}

\section{发布检查表}

\begin{enumerate}[leftmargin=2em]
  \item README 链接是否指向新仓库；
  \item bootstrap 是否指向新仓库；
  \item Windows 与 macOS installer 是否都包含 Claude + Codex；
  \item \texttt{.codex/.gitignore} 是否保护 auth/runtime；
  \item 真实运行 Windows installer 是否 exit 0；
  \item \texttt{codex --version} 是否可用；
  \item \texttt{code --list-extensions} 是否包含两个核心扩展；
  \item GitHub release 附件是否包含最新 bat/command。
\end{enumerate}

\section{Release 与 CI 边界}

当前仓库没有 \texttt{.github/workflows} 自动化发布配置。对外分发主要依赖手动创建 GitHub Release，并把最新的 \texttt{install-windows.bat} 与 \texttt{install-macos.command} 作为附件上传。双击入口本身并不包含完整安装逻辑，而是下载 GitHub 上对应 ref 的仓库压缩包再调用仓库内安装器。

\section{容易写错或理解错的点}

\begin{itemize}[leftmargin=2em]
  \item 不要说安装器会弹 API key 输入框；当前设计是安装后打开 cc-switch，由用户在 cc-switch 里配置 Claude provider。
  \item 不要说安装器会复制 Codex 认证文件；Codex 认证必须由 \texttt{codex login} 创建。
  \item 不要把 hooks 文件夹里的脚本总数等同于注册生命周期 hook 数；注册生命周期 hook 是 12 个，目录里还有辅助脚本。
  \item 不要把 \texttt{/mcp-status} 写成实时连接测试；它只检查配置，实时状态用 Claude Code 内置 \texttt{/mcp}。
  \item 不要说 Windows 安装 PowerShell 7；Windows 使用自带 PowerShell 5.1 即可，macOS 才会安装 PowerShell 7。
\end{itemize}

\chapter{附录：关键命令速查}

\section{安装}

\begin{lstlisting}[language={}]
# Windows
irm https://raw.githubusercontent.com/liujiarui0918/claude-code-codex-strongest/main/install/bootstrap.ps1 | iex

# macOS / Linux
curl -fsSL https://raw.githubusercontent.com/liujiarui0918/claude-code-codex-strongest/main/install/bootstrap.sh | bash
\end{lstlisting}

\section{验证}

\begin{lstlisting}[language={}]
claude --version
codex --version
codex doctor --summary --ascii --no-color
code --list-extensions
\end{lstlisting}

\section{登录}

\begin{lstlisting}[language={}]
# Claude Code: use cc-switch GUI
# Codex:
codex login
\end{lstlisting}

\backmatter
\printindex
\end{document}
